\documentclass[12pt]{article}
\usepackage[utf8]{inputenc}
\usepackage[T1]{fontenc}
\usepackage{hyperref}
\title{Wiraf\'ed Language Resources for K-12 Curriculum Development}
\author{Compiled by Claude Sonnet 5.5 (claude-sonnet-5-5) for Andr\'as Kornai}
\date{}
\begin{document}
\maketitle

\section{Introduction}
Wiraf\'ed (ISO 639-3: wir; Glottocode wira1264) is an extinct Tupian language of Brazil. Glottolog describes it as formerly spoken in the upper Rio Tapaj\'os area of Mato Grosso. Other names given are Tupi do Machado and Bocas Pretas. ASJP lists it as extinct since 1990. Sources disagree on whether it is a separate language or part of the Apiak\'a/Kawahib complex. Glottolog says it is hard to tell apart from Apiak\'a. The search found no teaching materials of any kind.

\section{Primary Educational Resources}
None found. No primers, textbooks or school programmes exist in the search results.

\section{Linguistic and Reference Materials}
Documentation is limited to old wordlists, attributed to Koch-Gr\"unberg, Loukotka and Nimuendaj\'u. The results gave no URLs for those wordlists.
\begin{itemize}
\item Glottolog entry for the language (ISO wir). \url{https://glottolog.org/resource/languoid/iso/wir}
\item ASJP database entry with a wordlist. \url{https://asjp.clld.org/languages/WIRAFED}
\item Numerals database page. It records traditional numbers 1--3, and notes that Portuguese numerals are now used. \url{https://lingweb.eva.mpg.de/channumerals/Wirafed.htm}
\item Wikipedia article on the related Apiak\'a language, which lists Wiraf\'ed as a dialect. \url{https://en.wikipedia.org/wiki/Apiak%C3%A1_language}
\end{itemize}

\section{Religious and Community Resources}
None found.

\section{Digital and Audio Resources}
None found. No audio or corpus was seen.

\section{Curriculum Development Context}
There are no speakers and almost no documentation, so a language curriculum is not feasible from what was found. Any heritage programme would have to rely on the few surviving wordlists and on comparison with related Kawahib and Apiak\'a materials. It should say plainly that the evidence base is minimal. Community descendants were not identified in the results.

\end{document}